\documentclass[14pt,a4paper]{extarticle}

% Margók beállítása
\usepackage[a4paper, left=3cm, right=2cm, top=3cm, bottom=3cm]{geometry}
\usepackage{float}

% Sortávolság 1.5
\usepackage{setspace}
\onehalfspacing

% Magyar nyelv támogatás
\usepackage[magyar]{babel}
\usepackage[utf8]{inputenc}
\usepackage[T1]{fontenc}
\usepackage{graphicx}

% Címlaphoz pozícionáláshoz
\usepackage{tikz}

\begin{document}

% --- CÍMLAP ---

\thispagestyle{empty}

% Egyetem, kar, tanszék (lap tetején, középre)
\vspace*{1cm}
\begin{center}
    \Large \textbf{Debreceni Egyetem} \\
    \large Informatika Kar \\
    \normalsize Információ Technológia Tanszék
\end{center}

\vspace{4cm}

% A dolgozat címe (felső harmad, középen)
\begin{center}
    \LARGE \textbf{ Közlekedési forgalom modellezése, elemzése és vizualizációja } \\
    \vspace{0.5cm}
\end{center}

\vspace{7cm}

% Alsó harmad, balra témavezető, jobbra szerző
\begin{minipage}[t]{0.45\textwidth}
    \raggedright
    \textbf{Témavezető:} \\    
Prof. Dr. Ispány Márton
 \\
    tanszékvezető, oktatási dékánhelyettes 
\end{minipage}
\hfill
\begin{minipage}[t]{0.45\textwidth}
    \raggedleft
    \textbf{Szerző:} \\
    Dáné Bence Dávid \\
    Adattudomány MSc
\end{minipage}

\vfill

\begin{center}
    Debrecen, 2026
\end{center}

\newpage


\setcounter{page}{1} % oldalszám 1-től a tartalomjegyzéknél
\tableofcontents

\newpage


\section{Bevezetés}
\subsection{Aktualitás}
A növekvő urbanizáció következményeként a városi közlekedés egy összetett és kritikus fontosságú infrastruktúrává vált. A népesség és járművek számának növekedése, a környezeti kihívások, valamint a lakosság mobilitási igényei nyomást gyakorolnak a városi úthálózatra és a közösségi közlekedés rendszereire. A valósághű szimulációs modellek és az adatalapú döntéshozatal lehetőséget nyújt különböző forgalmi, infrastrukturális és környezeti szcenáriók futtatására és elemzésére, anélkül, hogy hatalmas költségek halmozódnának fel.
\\
A városi közlekedés szimulációja kiemelten fontos, mert segítségével megérthetők, magyarázhatók és előre jelezhetők olyan jelenségek, mint a forgalmi torlódások vagy a közösségi közlekedés járműveinek késései és azok következményei. Egy jármű vagy esemény hatása nem triviális, és gyakran kihat a teljes hálózatra.
\\
Az adatalapú, szimuláción alapuló megközelítés nemcsak technikai kihívást jelent, hanem kiváló lehetőség egy rendszerszintű szemlélet kialakítására is. A városi közlekedés ugyanis nem csupán járművek mozgását jelenti, hanem összefonódik az emberek napi rutinjával, a gazdasággal, az időjárással, az eseményekkel, valamint a város tervezésének és irányításának kérdéseivel is.
\\
Ezért egy olyan szimulációs rendszer, amely a valósághoz közelítő módon modellezi a járművek – különösen a közösségi közlekedési járművek – viselkedését, késéseit, valamint a forgalmi eseményekre adott reakcióit, nemcsak tudományos szempontból érdekes, hanem valós alkalmazási lehetőségeket is kínál. Az ilyen modellek és szimulációk hozzájárulhatnak a döntéshozatal támogatásához, a várostervezés fejlesztéséhez, valamint a közlekedésirányítás optimalizálásához.
\subsection{Formális követelmények}
A diplomamunka célja egy adat által irányított, városi közlekedést modellező szimulációs rendszer létrehozása ami különösen hangsúlyozza a közösségi közlekedésjárművek, buszok és villamosok viselkedését, hálózati hatásait és késések terjedését.
A dolgozat középponti kérdései a következők: hogyan lehet különféle forrásból (pl. OpenStreetMap, GTFS) származó adatokat egyesíteni és feldolgozni úgy, hogy azok alapként szolgálhassanak egy valóságszerű közlekedési szimulációhoz? Miként képes egy közlekedési jármű késése - akár pár perces is -, láncreakciót okozni a teljes hálózaton? Hogyan lehet ezt a hatást modellezni és vizualizálni?
\\
A szimulációs motor alapja a diszkrét események által vezérelt időbeli történések (pl. megállások, balesetek, útlezárások). Különböző matematikai modellek kerülnek bemutatásra és alkalmazásra (jármű követés, sávváltás) az egyes járművekhez, illetve véletlenszerű események generálódnak Markov-láncokat használva.
\\
A program célja egy valós térképalapú vizualizációval rendelkező szimulációs eszköz, amely különböző szcenáriók (pl. csúcsforgalom, ünnepnap, baleset, időjárás) hatásait tudja bemutatni. A felhasználói felület lehetővé teszi a szimuláció futtatását különböző paraméterekkel, valamint az egyes járművek viselkedésének és késéseinek nyomon követését.
\section{Elméleti háttér}
A városi közlekedés egy összetett rendszer, az úthálózatokon mozgó járművek folyamatos interakciója, a közlekedési szabályok, emberi viselkedés és a környezeti hatások együttese megnehezíti a rendszer leírását és szimulálását. Ennek kezelésére szolgál a \textit{forgalomáramlás-elmélet} (\textit{traffic flow theory}), amely különböző szinteken értelmezi és modellezi az egyes járművek mozgását.
\\
\\
A forgalomáramlás (q) definiálható egy adott útszakaszon áthaladó járművek számaként egységnyi idő alatt, és összefügg a járműsűrűséggel (k) és az átlagsebességgel (v):
\\
$$q = k \cdot v$$
\begin{figure}[H]
    \centering
    \includegraphics[width=.45\linewidth]{t09-combined-fundamental-diagrams}
    \caption{Forgalomáramlás \cite{ref1}}
    \label{fig:flow-fund}
\end{figure}
\subsection{Forgalomáramlás modell}
\subsubsection{Makroszkopikus}
A makroszkopikus modellek a forgalom áramlását a mozgásban lévő folyadékokhoz vagy gázokhoz hasonlóan írják le.A dinamikai változók lokálisan aggregált mennyiségek, mint például a forgalomsűrűség $p(x,t)$, az áramlás $Q(x,t)$, az átlagsebesség $V(x,t)$.
Mivel az aggregáció lokális, ezek a mennyiségek általában térben és időben változnak, azaz dinamikus mezőknek felelnek meg. Így a makroszkopikus modellek képesek leírni az olyan kollektív jelenségeket, régiók torlódása vagy forgalmi hullámok. Továbbá, a makroszkopikus modellek hasznosak,
\begin{itemize}
    \item ha nem szükséges figyelembe venni olyan hatásokat, amelyeket makroszkopikusan nehéz leírni (pl. sávváltások, többféle vezető-jármű típus),
    \item ha csak makroszkopikus mennyiségek érdeklik az embert.
\end{itemize}
Egy makroszkopikus forgalomleírásban nem az egyedi járműveket írják le. Inkább az egyes útszakaszok aggregált változóit írják le. Azaz meg lehet határozni a sűrűséget, vagyis azt, hogy a járművek mennyire közel vannak egymáshoz térben, az áramlást, azaz az időegységenként egy referencia ponton áthaladó járművek számát és leírható a járművek átlagsebessége egy útszakaszon.
\begin{figure}[H]
    \centering
    \includegraphics[width=.5\linewidth]{Screenshot 2025-06-19 180354}
    \caption{Makroszkopikus modell. \cite{ref2}.}
    \label{fig:makro-model}
\end{figure}
\subsubsection{Mikroszkopikus}
A mikroszkopikus modellek, beleértve az autókövető modelleket, az egyéni járműveket írják, melyek együttesen alkotják a forgalom áramlását. Egy járműről teljes információt a trajektóriája ad, azaz a jármű pozíciójának minden időpillanatban történő meghatározása. A modellek a környező forgalomtól függően mutatják be az egyes vezetők reakcióját (gyorsítás, fékezés, sávváltás).  dinamikai változók a járművek pozíciói $x_{\alpha}(t)$, sebességei $v_{\alpha}(t)$ és gyorsulásai $\dot{v}_{\alpha}(t)$. A mikroszkopikus modellek alkalmasak a következő alkalmazásokra:
\begin{itemize}
	\item Az egyedi járművek forgalomra gyakorolt hatásának modellezése,
	\item Olyan helyzetek, amelyekben a forgalom heterogenitása fontos szerepet játszik, pl. a sebességkorlátozások vagy teherautók előzési tilalmának hatásainak szimulálása,
	\item Különböző forgalmi résztvevők (autók, teherautók, buszok, kerékpárosok, gyalogosok stb.) közötti interakciók vizualizálása,
	\item Az emberi vezetési viselkedés leírása, beleértve a becslési hibákat, reakcióidőket, figyelmetlenséget és előrelátást.
\end{itemize}
\begin{figure}[H]
	\centering
	\includegraphics[width=.55\linewidth]{micro-fig}
	\caption{Mikroszkopikus modell. \cite{ref2}}
	\label{fig:micro-fig}
\end{figure}
A megfelelő modell, vagy modellek kiválasztásakor (mezoszkopikus modell, amely ötvözi a mikroszkopikus és makroszkopikus megközelítéseket) az alábbi tényezők is szerepet játszanak:
\begin{itemize}
	\item \textit{Heurisztikus} és \textit{Alapelven nyugvó} modellek.  A \textit{heurisztikus modellek} egyszerű matematikai megközelítést használnak, ahol az együtthatók a modellparaméterek szerepét töltik be. Ezeket az adatokat például regressziós technikákkal illesztik, és általában nincs intuitív jelentésük. Ezzel szemben az \textit{alapelveken nyugvó modelleket} bizonyos szükségességekből vezetik le. Autókövető modellek esetében ez lehet egy olyan vezetési viselkedés, amelyet a sebességre, gyorsulásra, lassulásra, időbeli követési távolságra és minimális távolságra vonatkozó kívánt értékek határoznak meg. Ideális esetben ezen szükségességek mindegyike egy modellparaméterben tükröződik, amelynek értéke ezáltal intuitív jelentéssel bír.
	\item \textit{Sztochasztikus} és \textit{Determinisztikus} modellek. A véletlen elemek felhasználhatók a forgalomáramlás olyan aspektusainak leírására, amelyek ismeretlenek, mérhetetlenek, modellezhetetlenek, vagy „valóban” véletlenszerűek. Míg a véletlen elemeket nem tartalmazó modelleket \textit{determinisztikus} modelleknek nevezzük, addig azokat, amelyek véletlen elemeket (sztochasztikus tagokat) tartalmaznak, \textit{sztochasztikus} modelleknek hívjuk.A véletlenszerűség a modell különböző pontjain előfordulhat:
	\begin{itemize}
		\item gyorsulási zaj, amely a emberi viselkedés kiszámíthatatlanságát modellezi
		\item exogén zaj, amelyet a bemeneti adatokhoz adnak hozzá, és az emberi érzékelési és becslési hibáinak modellezésére szolgál 
	\end{itemize}
	\item \textit{Egysávos} és \textit{Többsávos} modellek. Ha a forgalomáramlási modell több sávot és a köztük lévő sávváltásokat is leírja, akkor két fő komponensből áll: a longitudinális dinamikából (gyorsulási modell) és a laterális dinamikából (sávváltási modell). Egyes modellek eredendően tartalmazzák a laterális dinamikát, míg a tisztán longitudinális modellek kiegészíthetők egy megfelelő sávváltási modellel.
\end{itemize}
A motorizált forgalomáramlás mellett a fentebb tárgyalt modellkeretekbe beilleszthető a nem motorizált forgalom dinamikája is. A nem motorizált forgalom magában foglalja a gyalogos, kerékpáros és vegyes forgalmat.  A járművekkel ellentétben a gyalogosok általában szabadon mozoghatnak két térdimenzióban, azaz két térbeli koordinátájuk van (x és y). A kívánt (járási) sebesség mellett minden gyalogosnak van egy kívánt iránya is. Következésképpen a kívánt sebesség vektorális mennyiség.
\subsection{Autókövető modellek}
Az autókövető modellek a mikroszkopikus forgalomáramlási modellek legfontosabb képviselői. Ezek egyedi vezető-jármű egységek szemszögéből írják le a forgalom dinamikáját. Szigorú értelemben az autókövető modellek csak a vezető viselkedését írják le más járművekkel való interakciók jelenlétében, míg a szabad forgalomáramlást külön modell írja le.
Egy autókövető modell akkor teljes, ha képes leírni minden helyzetet, beleértve a szabad forgalomban való gyorsítást és utazósebességet, a többi jármű követését álló és nem álló helyzetekben, valamint a lassú vagy álló járművek, és a piros lámpák megközelítését.
A modellezett cselekvésektől függően megkülönböztetünk \textit{gyorsulási modelleket} a longitudinális mozgáshoz, \textit{sávváltási modelleket} a laterális mozgáshoz, és \textit{döntési modelleket} más diszkrét választási helyzetekhez, mint például egy elsőbbségi útra való behajtás.
\subsection{Lane-Changing Models}
\subsection{Public Transport Models}
\subsection{Discrete-Event Simulation (DES)}
\subsection{Modeling Random Events}
\subsection{Summary:How these models interact in your system, Which agents use which model,
	How this sets the stage for your architecture}

\section{Összefoglalás}
Röviden, tömören és világosan ismertetni kell a fontosabb megállapításokat és következtetéseket.

\section{Irodalomjegyzék}
\begin{thebibliography}{9}
\bibitem{ref1} Dr. Tom V. Mathew. Fundamental relation of traffic flow. 
\bibitem{ref2} Treiber, M. and Kesting, A. (2013) Traffic Flow Dynamics: Data, Models and Simulation. Springer.
\bibitem{ref3} Victor L. Knoop, (2018) Traffic Flow Theory: An Introduction with Exercises. TU Delft Open.
\end{thebibliography}

\end{document}
